\section{Bitter Lesson: no ir en contra de la mayor variable de la era de los modelos}

La sección anterior trató la base del producto; esta sección aborda el primer gran giro en el emprendimiento de aplicaciones de IA. Poco después de dejar Moonshot para emprender, Ming Chaoping (明超平) se topó con la bitter lesson: cuando la capacidad de los modelos es la mayor variable de la época, si una empresa de aplicaciones añade por instinto demasiado andamiaje, plantillas y ornamentos al producto, puede terminar alejándose de la dirección en la que avanzan los modelos. Detener el producto antes de lanzarlo no se debió a falta de trabajo, sino a que se dieron cuenta de que habían apostado por la variable equivocada.

\begin{figure}[H]
\centering
\includegraphics[width=0.96\textwidth]{bitter-lesson-pivot.png}
\caption{Giro de la Bitter Lesson: querer ornamentar y añadir andamiaje puede alejar el producto de la mayor variable de la época. Diagrama conceptual propio, elaborado a partir de la conversación de 01:05:05--01:13:18.}
\end{figure}

\begin{knowledgebox}{Lectura de la figura: saber detenerse también es una capacidad de producto}
La figura va de Scaffold a Launch?, Wrong variable, Stop y Rethink. Lo esencial aquí no es «no hacer nada», sino que la empresa de aplicaciones debe juzgar qué experiencias debe asumir el producto y qué capacidades conviene dejar que surjan de forma natural en el modelo. Confundir el andamiaje con una ventaja defensiva hace que el producto se vuelva cada vez más pesado.
\end{knowledgebox}

\subsection{El auge y la caída repentina de Noisee}

Esta subsección completa un giro importante de la etapa en Moonshot. En la entrevista se menciona el auge y la caída repentina de Noisee, un producto del extranjero, que le mostró a Ming Chaoping dos rasgos de la ventana de oportunidad de los productos de IA: el despegue puede llegar muy rápido, pero saber si el producto realmente sigue las capacidades del modelo y los escenarios de uso exige un juicio más sereno. Un producto puede recibir retroalimentación positiva de corto plazo en difusión, pero si la variable de fondo es incorrecta, seguir ornamentándolo solo aleja al equipo de la línea principal de la época.

Esta experiencia y la bitter lesson se explican mutuamente. Emprender en aplicaciones de IA no consiste en «hacer un demo vistoso y luego pagar por tráfico», sino en juzgar qué tipo de capacidad están liberando los modelos y en qué entorno los usuarios pueden convertir esa capacidad más rápido en obras o en productividad. Noisee le dio una percepción directa de la frontera; YouWare, en cambio, fue la elección de un nuevo punto de entrada, más cercano a la creación con coding y al entorno de publicación.

\begin{warningbox}{La retroalimentación positiva en difusión no significa que el rumbo del producto sea correcto}
Los productos de IA logran difusión con facilidad gracias a la novedad, pero la difusión no sustituye la retención, la creación, la reutilización ni el ecosistema. Los giros de Ming Chaoping en Moonshot y en las primeras etapas de su emprendimiento consistieron precisamente en distinguir entre «volverse popular» y «tener el rumbo correcto».
\end{warningbox}

\subsection{Velocidad de consumo de tokens y per token valuation}

La revelación que tuvo tras una noche de insomnio fue que uno de los indicadores clave de la era de la IA podría ser la velocidad de consumo de tokens. Los usuarios no gastan tokens por «gastar tokens», sino para que la inteligencia se convierta más rápido en resultados. Un buen escenario de Agent debería completar tareas más grandes con menos tokens, o hacer que los tokens se conviertan con mayor eficiencia en obras, sitios web, juegos, procesos y productividad.

\begin{figure}[H]
\centering
\includegraphics[width=0.96\textwidth]{token-consumption-speed.png}
\caption{Velocidad de consumo de tokens: en la era de la IA, los indicadores de producto pasan de clics/tiempo de uso a un consumo más eficiente de inteligencia. Diagrama conceptual propio, elaborado a partir de la conversación de 01:13:18--01:16:33 y 02:03:33--02:14:50.}
\end{figure}

\begin{knowledgebox}{Lectura de la figura: el token es consumo de inteligencia, no solo un costo}
En la figura, User intent entra al Container, luego Agent action consume Tokens y finalmente estos se convierten en Value. Al leer la figura, no hay que ver el token solo como costo de API; para una empresa de aplicaciones, el token es la unidad con la que se moviliza la inteligencia, y lo decisivo es cuánto valor para el usuario crea cada token.
\end{knowledgebox}

\[
\mathrm{PTV}=\frac{\mathrm{User\ Value}}{\mathrm{Tokens\ Consumed}}
\]

Donde $\mathrm{PTV}$ denota la per token valuation; $\mathrm{User\ Value}$ es la obra, la eficiencia, el ingreso o la experiencia que el usuario obtiene realmente; y $\mathrm{Tokens\ Consumed}$ son los tokens del modelo consumidos para completar la tarea. Esta expresión no es una fórmula financiera, sino una herramienta para pensar el producto: si se consumen muchos tokens y solo se produce una conversación casual, la densidad de valor es baja; si con pocos tokens se genera un sitio web o un juego que se puede compartir, con el que se puede interactuar y que se puede seguir editando, la densidad de valor es alta.

\begin{warningbox}{No interpretar el consumo de tokens como «cuanto más, mejor»}
La velocidad de consumo importa, pero no debe entenderse burdamente como quemar tokens. El objetivo real es convertir la inteligencia, de forma más rápida y eficiente, en valor perceptible para el usuario. Las conversaciones largas sin propósito, el ensayo y error repetido y el razonamiento en vacío solo elevan el costo y no forman una ventaja defensiva para el producto.
\end{warningbox}

\subsection{Resumen de la sección}

La primera lección del emprendimiento en aplicaciones de IA es seguir la mayor variable de la era de los modelos, en lugar de envolver al modelo en una capa compleja por hábitos de productos anteriores. Los indicadores clave de YouWare pasan de los clics/tiempo de uso tradicionales a cómo la inteligencia se encapsula en un contenedor, se consume con eficiencia y se convierte en obras.
